% Define o tipo de documento como "article", com fonte de 12pt
% A opção 'nodisplayskipstretch' evita que o LaTeX altere o espaçamento ao redor de fórmulas exibidas
\documentclass[12pt,nodisplayskipstretch]{article}

% Importa o arquivo com as configurações gerais do projeto (pacotes, margens, fonte etc.)
\usepackage{config/commands}

% Importa o arquivo de configurações adicionais (talvez pacotes específicos ou definições de estilo)
\usepackage{config/config}

\begin{document} % Início do conteúdo do documento

% ----------------------------
% SEÇÕES DO DOCUMENTO
% ----------------------------

% Importa a capa do trabalho
\import{sections/}{0-capa.tex}

% Importa a seção de identificação do projeto (pode conter título, autor, orientador, etc.)
\import{sections/}{1-identificacao.tex}

% Importa a introdução do projeto, onde o tema é apresentado
\import{sections/}{2-introducao.tex}

% Importa a seção de objetivos do projeto, onde são definidos os objetivos gerais e específicos
\import{sections/}{3-objetivos.tex}

% Importa a justificativa do projeto, onde são apresentados os motivos e a relevância do trabalho
\import{sections/}{4-justificativa.tex}

% Importa a seção de metodologia, onde são descritos os métodos que serão utilizados no projeto
\import{sections/}{5-metodologia.tex}

% Importa a seção de equipamento e material necessário para o desenvolvimento do projeto
\import{sections/}{6-equipamento-material.tex}

% Importa o cronograma de execução do projeto, com as etapas e os prazos
\import{sections/}{7-cronograma-de-execucao.tex}

% Importa as referências bibliográficas utilizadas no projeto
\import{sections/}{8-referencias.tex}

\end{document} % Fim do conteúdo do documento
